% Appendix C: Tower Calculus Proofs (Chapter 5)
\label{sec:appendix-tower}

% ============================================================
% Proofs from foundational_lemmas.tex
% ============================================================

\begin{proof}[Proof of Lemma~\ref{lem:block-splitting}]
\label{proof:lem:block-splitting}
Apply Lemma~\ref{lem:resolution-identity} to $B\CCPi{h}x$:
\[
\CCPi{h}AB\CCPi{h}
=
\CCPi{h}A(\CCPi{h} + \CCPibar{h})B\CCPi{h}
=
\CCPi{h}A\CCPi{h}B\CCPi{h} + \CCPi{h}A\CCPibar{h}B\CCPi{h}.
\]
The first term is the observed composition and the second is the cross-shadow correction.
\end{proof}

% ============================================================
% Proofs from hft.tex
% ============================================================

\begin{proof}[Proof of Lemma~\ref{lem:silent-leak}]
\label{proof:lem:silent-leak}
Take the rebuilt finite state space and let $\Pi$ project onto the first coordinate:
\[
\Pi(x_0,x_1,x_2,x_3) = (x_0,0,0,0).
\]
Let $d$ send the first coordinate into the second and annihilate the others:
\[
d(x_0,x_1,x_2,x_3) = (0,x_0,0,0).
\]
Then $d^2 = 0$. The observed operator $d_\Pi = \Pi d \Pi$ vanishes, so $d_\Pi^2 = 0$.
But $L(1,0,0,0) = (0,1,0,0) \neq 0$, while the return map vanishes on that shadow image, so $\delta = RL = 0$.
\end{proof}

% ============================================================
% Proofs from tower_calculus.tex
% ============================================================

\begin{proof}[Proof of Lemma~\ref{lem:tower-derivative-block}]
\label{proof:lem:tower-derivative-block}
On the $(h,k)$ block,
\[
\Delta_h \CCIndexOp A \Delta_k = h\, \CCBlock{A}{h}{k},
\qquad
\Delta_h A \CCIndexOp \Delta_k = k\, \CCBlock{A}{h}{k},
\]
so the commutator contributes the factor $(h-k)$. The Leibniz rule is the
usual commutator identity
\[
[\CCIndexOp, AB] = [\CCIndexOp, A]B + A[\CCIndexOp, B].
\]
\end{proof}

\begin{proof}[Proof of Proposition~\ref{prop:tower-order-arithmetic}]
\label{proof:prop:tower-order-arithmetic}
Statement (i) is immediate from the definition of block support.

For (ii), if $|j-k|>m+n$, then for every intermediate band index $\ell$ one has
either $|j-\ell|>m$ or $|\ell-k|>n$. Hence every term in the block expansion
\[
\CCDelta{j}AB\CCDelta{k}
=
\sum_{\ell}\CCDelta{j}A\CCDelta{\ell}B\CCDelta{k}
\]
vanishes, so $\CCDelta{j}AB\CCDelta{k}=0$.

For (iii), Lemma~\ref{lem:tower-derivative-block} gives
\[
\CCBlock{\CCDer{A}}{j}{k}=(j-k)\,\CCBlock{A}{j}{k}.
\]
Thus the block support of $\CCDer{A}$ is exactly the block support of $A$.
\end{proof}

\begin{proof}[Proof of Proposition~\ref{prop:transport-recursion-compatibility}]
\label{proof:prop:transport-recursion-compatibility}
For $g\in G$ and $x\in F^{\bbZ}$,
\[
(\mathcal R_B\mathcal U(g)x)_n
=
U(g)x_{n+1}-BU(g)x_n+U(g)x_{n-1}
=
U(g)(x_{n+1}-Bx_n+x_{n-1})
=
(\mathcal U(g)\mathcal R_Bx)_n,
\]
which proves (a). The same calculation with $J_F$ in place of $U(g)$ gives
\[
(\mathcal R_B\mathcal Jx)_n
=
J_Fx_{n+1}-BJ_Fx_n+J_Fx_{n-1}
=
J_F(x_{n+1}-Bx_n+x_{n-1})
=
(\mathcal J\mathcal R_Bx)_n,
\]
proving (b).

For (c), let $x,y$ be finitely supported sequences. Using the sum inner
product and shifting indices in the finite sums,
\begin{align*}
\langle \mathcal R_Bx,y\rangle
&=
\sum_n \langle x_{n+1},y_n\rangle
-\sum_n \langle Bx_n,y_n\rangle
+\sum_n \langle x_{n-1},y_n\rangle \\
&=
\sum_n \langle x_n,y_{n-1}\rangle
-\sum_n \langle x_n,B^*y_n\rangle
+\sum_n \langle x_n,y_{n+1}\rangle \\
&=
\sum_n \langle x_n,y_{n-1}-By_n+y_{n+1}\rangle \\
&=
\langle x,\mathcal R_By\rangle,
\end{align*}
because $B=B^*$.
\end{proof}

\begin{proof}[Proof of Proposition~\ref{prop:repeat-cycle-monodromy-factorization}]
\label{proof:prop:repeat-cycle-monodromy-factorization}
By definition, the monodromy of \(W^{\times m}\) is the product of the same
block \(T_r\cdots T_1\) repeated \(m\) times. Associativity of composition gives
\[
\mathcal M(W^{\times m})
=
(T_r\cdots T_1)\cdots(T_r\cdots T_1)
=
\mathcal M(W)^m.
\]
The final statement is exactly this factorization.
\end{proof}

\begin{proof}[Proof of Proposition~\ref{prop:cyclic-refinement-prime-step}]
\label{proof:prop:cyclic-refinement-prime-step}
Since \(H\) is a subgroup of a cyclic group, it is itself cyclic. Write
\[
H\cong C_m
\]
for some \(m\). Subgroups of a cyclic group correspond to divisors of its
order. If \(H'\subsetneq H\) is maximal, then there is no intermediate
subgroup, equivalently no divisor strictly between \(|H'|\) and \(|H|\).
Therefore the quotient \(|H|/|H'|\) has no nontrivial factor and must be
prime. Since \([H:H']=|H|/|H'|\), the index is prime.
\end{proof}

\begin{proof}[Proof of Theorem~\ref{thm:cyclic-refinement-prime-invariant}]
\label{proof:thm:cyclic-refinement-prime-invariant}
Let
\[
C_k=H_0\supset H_1\supset\cdots\supset H_r=\{0\}
\]
be a maximal refinement chain. Taking orders gives
\[
k=|C_k|=\prod_{i=0}^{r-1}\frac{|H_i|}{|H_{i+1}|}
=
\prod_{i=0}^{r-1}[H_i:H_{i+1}].
\]
By Proposition~\ref{prop:cyclic-refinement-prime-step}, each factor is prime.
Because \(C_k\) is cyclic, for every divisor \(d\mid k\) there is a unique
subgroup of order \(d\). Hence a maximal refinement chain is exactly a strict
descending divisor chain with no intermediate divisor steps, so its step
indices are prime and their product is \(k\). The uniqueness of the multiset of
prime factors of \(k\) therefore implies that every maximal cyclic refinement
chain has the same multiset of prime indices, up to permutation.
\end{proof}

\begin{proof}[Proof of Theorem~\ref{thm:isotropic-quadratic-transport}]
\label{proof:thm:isotropic-quadratic-transport}
Because $B$ commutes with the orthogonal action $U(g)$, the operator $B^*B$
also commutes with every $U(g)$. It is self-adjoint by construction. Since the
representation of $K$ on $F$ is irreducible, Corollary~\ref{cor:self-adjoint-commutant-scalarization}
forces every self-adjoint operator commuting with the action to be a scalar
multiple of the identity. Thus
\[
B^*B=\kappa(B)\,I_F
\]
for a unique real scalar $\kappa(B)$. Positivity of $B^*B$ implies
$\kappa(B)\ge 0$. The energy formula is then immediate.
\end{proof}

\begin{proof}[Proof of Theorem~\ref{thm:sphere-transitive-implies-irreducible}]
\label{proof:thm:sphere-transitive-implies-irreducible}
Suppose $W\subseteq F$ is a nonzero proper $K$-invariant subspace. Choose a
unit vector $w\in W$. Since $W$ is proper, there exists a unit vector
$v\in W^{\perp}$. By sphere transitivity, there exists $g\in K$ such that
\[
U(g)w=v.
\]
But $W$ is $K$-invariant, so $U(g)w\in W$, hence $v\in W\cap W^{\perp}$, which
forces $v=0$. This contradicts the choice of $v$ as a unit vector. Therefore
no nonzero proper invariant subspace exists, and $F$ is irreducible.
\end{proof}

\begin{proof}[Proof of Theorem~\ref{thm:metric-isotropic-quadratic-scalar}]
\label{proof:thm:metric-isotropic-quadratic-scalar}
Choose an orthonormal basis $\{e_i\}_{i=1}^m$ of $F$. By metric isotropy,
\[
\langle Qe_i,e_i\rangle=\kappa(Q)
\qquad(i=1,\dots,m).
\]
Thus every diagonal coefficient of $Q$ in this basis equals $\kappa(Q)$.

For $i\neq j$, the unit vector
\[
u_{ij}^{+}:=\frac{e_i+e_j}{\sqrt2}
\]
satisfies
\[
\kappa(Q)=\langle Qu_{ij}^{+},u_{ij}^{+}\rangle
=
\frac12\langle Qe_i,e_i\rangle
\frac12\langle Qe_j,e_j\rangle
\langle Qe_i,e_j\rangle
=
\kappa(Q)+\langle Qe_i,e_j\rangle.
\]
Hence
\[
\langle Qe_i,e_j\rangle=0
\qquad(i\neq j).
\]
So the matrix of $Q$ in the orthonormal basis $\{e_i\}$ is diagonal with every
diagonal entry equal to $\kappa(Q)$, which means
\[
Q=\kappa(Q)\,I_F.
\]
\end{proof}

\begin{proof}[Proof of Theorem~\ref{thm:transitive-direction-isotropy}]
\label{proof:thm:transitive-direction-isotropy}
For $g\in K$,
\[
U(g)QU(g)^*
=
\sum_{\alpha\in\mathcal D} q_{\alpha}\,U(g)\Pi_{\alpha}U(g)^*
=
\sum_{\alpha\in\mathcal D} q_{\alpha}\,\Pi_{g\cdot\alpha}.
\]
Because $Q$ commutes with the action, this must equal
\[
Q=\sum_{\alpha\in\mathcal D} q_{\alpha}\Pi_{\alpha}.
\]
Since the group acts transitively on $\mathcal D$, the coefficients are
constant on the orbit and hence there exists a real number $q$ such that
\[
q_{\alpha}=q
\qquad(\alpha\in\mathcal D).
\]
Therefore
\[
Q=q\sum_{\alpha\in\mathcal D}\Pi_{\alpha}=qc\,I_F.
\]
Setting $\kappa(Q):=qc$ gives the claim, and uniqueness is immediate.
\end{proof}

\begin{proof}[Proof of Theorem~\ref{thm:quadratic-transport-rigidity}]
\label{proof:thm:quadratic-transport-rigidity}
If $B$ carries data of type (i), the conclusion is exactly
Theorem~\ref{thm:isotropic-quadratic-transport}. If it carries data of type
(ii), the conclusion is exactly
Corollary~\ref{cor:sphere-transitive-quadratic-transport}. If it carries data
of type (iii), the conclusion is exactly
Corollary~\ref{cor:full-orthogonal-quadratic-transport}. If it carries data of
type (iv), the conclusion is exactly
Corollary~\ref{cor:metric-isotropic-transport}. Finally, if it carries data of
type (v), the conclusion is exactly
Corollary~\ref{cor:direction-frame-quadratic-coefficient}. In every case
$B^*B$ is positive semidefinite, so the scalar $\kappa(B)$ is nonnegative, and
uniqueness is immediate.
\end{proof}

\begin{proof}[Proof of Theorem~\ref{thm:orbit-generated-direction-frame}]
\label{proof:thm:orbit-generated-direction-frame}
By construction, each $\Pi_{\alpha}$ is a rank-one orthogonal projection and
the group acts transitively on the orbit $\mathcal D$. Define
\[
S:=\sum_{\alpha\in\mathcal D}\Pi_{\alpha}.
\]
Since the family is a single orbit, $S$ commutes with the orthogonal action of
$K$. It is also self-adjoint. Because $(F,U)$ is an isotropic transport
channel, the representation of $K$ on $F$ is irreducible. Therefore
Corollary~\ref{cor:self-adjoint-commutant-scalarization} forces
\[
S=c\,I_F
\]
for some real scalar $c$. Taking traces gives
\[
|\mathcal D|
=
\sum_{\alpha\in\mathcal D}\operatorname{Tr}(\Pi_{\alpha})
=
\operatorname{Tr}(S)
=
c\,\dim F,
\]
because each rank-one projection has trace $1$. Hence
\[
c=\frac{|\mathcal D|}{\dim F}.
\]
This proves both statements.
\end{proof}

\begin{proof}[Proof of Theorem~\ref{thm:ftc-tower}]
\label{proof:thm:ftc-tower}
If $h=k$, the right-hand side has block $\CCBlock{A}{h}{h}$. If $h \neq k$,
Lemma~\ref{lem:tower-derivative-block} gives
\[
\CCBlock{(\CCInt{\CCDer{A}})}{h}{k}
=
\frac{1}{h-k}\, \CCBlock{\CCDer{A}}{h}{k}
=
\frac{1}{h-k}(h-k)\, \CCBlock{A}{h}{k}
=
\CCBlock{A}{h}{k}.
\]
So every block agrees.
\end{proof}

\begin{proof}[Proof of Proposition~\ref{prop:tower-ibp}]
\label{proof:prop:tower-ibp}
Using the block formula from Lemma~\ref{lem:tower-derivative-block},
\[
\Tr(\CCDer{A}\, B)
=
\sum_{h,k} (h-k)\, \Tr(\CCBlock{A}{h}{k}\CCBlock{B}{k}{h}),
\]
while
\[
\Tr(A\, \CCDer{B})
=
\sum_{h,k} (k-h)\, \Tr(\CCBlock{A}{h}{k}\CCBlock{B}{k}{h}).
\]
The two sums are negatives of each other.
\end{proof}

% ============================================================
% Proofs from defect_rules.tex
% ============================================================

\begin{proof}[Proof of Proposition~\ref{prop:product-rule}]
\label{proof:prop:product-rule}
Decompose $x = \CCPi{h} x + \CCPibar{h} x$ and $y = \CCPi{h} y + \CCPibar{h} y$.
By bilinearity,
\begin{align*}
B(x,y) &= B(\CCPi{h} x + \CCPibar{h} x,\, \CCPi{h} y + \CCPibar{h} y) \\
&= B(\CCPi{h} x, \CCPi{h} y) + B(\CCPi{h} x, \CCPibar{h} y) + B(\CCPibar{h} x, \CCPi{h} y) + B(\CCPibar{h} x, \CCPibar{h} y).
\end{align*}
Apply $\CCPi{h}$ and use the closure condition $\CCPi{h} B(\CCPi{h} x, \CCPi{h} y) = B(\CCPi{h} x, \CCPi{h} y)$:
\begin{align*}
\CCTauBi{h}{B}{x}{y} &= \CCPi{h} B(x,y) - B(\CCPi{h} x, \CCPi{h} y) \\
&= B(\CCPi{h} x, \CCPi{h} y) + \CCPi{h} B(\CCPi{h} x, \CCPibar{h} y) + \CCPi{h} B(\CCPibar{h} x, \CCPi{h} y) + \CCPi{h} B(\CCPibar{h} x, \CCPibar{h} y) - B(\CCPi{h} x, \CCPi{h} y)\\
&= \CCPi{h} B(\CCPi{h} x, \CCPibar{h} y) + \CCPi{h} B(\CCPibar{h} x, \CCPi{h} y) + \CCPi{h} B(\CCPibar{h} x, \CCPibar{h} y). \qedhere
\end{align*}
\end{proof}

% ============================================================
% Proofs from constraint_compiler.tex
% ============================================================

\begin{proof}[Proof of Theorem~\ref{thm:compiler-contract}]
\label{proof:thm:compiler-contract}
This is a schema theorem asserting that the calculus provides a compositional route to a well-posed completion constraint.
\begin{itemize}[itemsep=0.1em]
\item Defect propagation (Section~\ref{sec:defect-rules}) composes via chain/product rules, yielding a derived defect contribution.
\item Elimination (Section~\ref{sec:resolvent-defect}) produces residual terms (Schur corrections) on the observable space.
\item These contributions aggregate into a ``required lower bound'' $D$ for stability.
\item Admissibility (Definition~\ref{def:admissibility}) selects completions $Q$ satisfying $D + Q \succeq 0$ that are minimal in PSD order.
\item Minimality in a partial order either yields a unique minimum or a minimal face (non-uniqueness certificate).
\end{itemize}
The theorem does not claim uniqueness of $D$; it states that once a bound is fixed, the completion constraint and certificate mechanism are canonical.
\end{proof}

\begin{proof}[Proof of Proposition~\ref{prop:operator-norm-selector-exists}]
\label{proof:prop:operator-norm-selector-exists}
By Definition~\ref{def:minimal-scalar-repair}, the operator
\[
\nu(D)I_{\mathcal K}
\]
is admissible, so the admissible set
\[
\mathcal A
:=
\{Q\in\mathcal Q: D+Q\succeq 0\}
\]
is nonempty. It is closed because the positive semidefinite cone is closed and
$Q\mapsto D+Q$ is continuous. Since $\mathcal Q$ is finite-dimensional, any
norm-minimizing sequence in $\mathcal A$ is bounded and therefore has a
convergent subsequence whose limit still lies in $\mathcal A$. Continuity of
the operator norm gives a minimizer.
\end{proof}

\begin{proof}[Proof of Theorem~\ref{thm:minimal-selectors-vanish}]
\label{proof:thm:minimal-selectors-vanish}
By Definition~\ref{def:minimal-scalar-repair}, the scalar correction
\[
\nu(D_h)I_{\CCPi{h}\cH}
\]
is admissible. Since $C_h^\ast$ minimizes operator norm among admissible
corrections,
\[
\|C_h^\ast\|
\le
\|\nu(D_h)I_{\CCPi{h}\cH}\|
=
\nu(D_h),
\]
proving (i). For $x\in\cH$,
\[
\|C_h^\ast \CCPi{h}x\|
\le
\|C_h^\ast\|\,\|\CCPi{h}x\|
\le
\nu(D_h)\,\|\CCPi{h}x\|.
\]
Because $\CCPi{h}$ is an orthogonal projection, $\|\CCPi{h}x\|\le \|x\|$, and
since $\nu(D_h)\to 0$, statement (ii) follows. Statement (iii) is exactly
Theorem~\ref{thm:selected-schur-recovery-general}.
\end{proof}

\begin{proof}[Proof of Theorem~\ref{thm:normal-form}]
\label{proof:thm:normal-form}
Each isotypic projector $P_\rho$ is constructed from the group action via $P_\rho = \frac{d_\rho}{|G|} \sum_{g \in G} \overline{\chi_\rho(g)}\, U(g)$, so $P_\rho$ commutes with every $U(g)$.
Since $A$ commutes with every $U(g)$ by hypothesis, $A$ commutes with $P_\rho$.

For distinct irreps $\rho \ne \rho'$, orthogonality of characters gives $P_\rho P_{\rho'} = 0$.
Thus:
\[
P_\rho \, A \, P_{\rho'} = P_\rho \, P_{\rho'} \, A = 0 \cdot A = 0.
\]
Since $\sum_\rho P_\rho = I$ on $\cH_{\mathrm{obs}}$, we have $A = \sum_{\rho, \rho'} P_\rho A P_{\rho'} = \sum_\rho P_\rho A P_\rho$.
\end{proof}

\begin{proof}[Proof of Theorem~\ref{thm:equivariant-coupling-decomposition}]
\label{proof:thm:equivariant-coupling-decomposition}
Let $P_\rho^X$ and $P_\rho^Y$ be the isotypic projectors for $X$ and $Y$.
Because $T$ intertwines the $G$-actions, it commutes with the projector
construction on both sides:
\[
P_\rho^Y T = T P_\rho^X
\qquad\text{for every }\rho.
\]
If $\rho\neq \rho'$, then orthogonality of isotypic projectors gives
\[
P_\rho^Y T P_{\rho'}^X = P_\rho^Y P_{\rho'}^Y T = 0,
\]
so $T$ is block-diagonal across isotypic channels:
\[
T=\sum_\rho P_\rho^Y T P_\rho^X.
\]
On the $\rho$-channel, Schur's lemma gives the standard normal form
\[
P_\rho^Y T P_\rho^X \cong I_{V_\rho}\otimes T_\rho
\]
for a unique linear map $T_\rho:M_\rho^X\to M_\rho^Y$. This proves the stated
decomposition.

For the dimension formula, each block contributes
\[
\dim \mathrm{Hom}_G(V_\rho\otimes M_\rho^X,\;V_\rho\otimes M_\rho^Y)
=
(\dim M_\rho^X)(\dim M_\rho^Y)\,\dim \mathrm{End}_G(V_\rho),
\]
and the block decomposition is a direct sum across $\rho$. Summing over the
appearing channels gives the stated formula.
\end{proof}

\begin{proof}[Proof of Theorem~\ref{thm:multiplicity-free-closures-scalar}]
\label{proof:thm:multiplicity-free-closures-scalar}
By Corollary~\ref{cor:isotypic-block}, every $G$-equivariant operator on the
$\rho$-isotypic component has the form
\[
P_{\rho}AP_{\rho}\cong I_{V_{\rho}}\otimes B_{\rho}
\]
for some operator $B_{\rho}$ on the multiplicity space $M_{\rho}$. If the
representation is multiplicity-free, then $\dim M_{\rho}=1$ for every appearing
$\rho$. Hence each $B_{\rho}$ is multiplication by a real scalar
$\lambda_{\rho}$ because $A$ is self-adjoint. Therefore $A$ is scalar on every
channel, which is exactly Definition~\ref{def:symmetry-scalar-closure}.
\end{proof}

\begin{proof}[Proof of Theorem~\ref{thm:orthogonal-symmetric-tensor-scalarization}]
\label{proof:thm:orthogonal-symmetric-tensor-scalarization}
Because $QIQ^\top=I$ for every $Q\in O(d)$, the equivariance relation gives
\[
L(I)=Q\,L(I)\,Q^\top
\qquad\text{for every }Q\in O(d).
\]
Thus $L(I)$ is fixed by every orthogonal conjugation, so there is a unique
scalar $\beta$ with
\[
L(I)=\beta I.
\]

For $i<j$, let
\[
F_{ij}:=E_{ij}+E_{ji},
\qquad
D_{ij}:=E_{ii}-E_{jj},
\]
where $E_{ab}$ are the standard matrix units. First consider $F_{12}$. Let
$R_1=\operatorname{diag}(-1,1,\dots,1)$ and
$R_2=\operatorname{diag}(1,-1,1,\dots,1)$. Then
\[
R_1F_{12}R_1^\top=-F_{12},
\qquad
R_2F_{12}R_2^\top=-F_{12}.
\]
Equivariance therefore yields
\[
R_1L(F_{12})R_1^\top=-L(F_{12}),
\qquad
R_2L(F_{12})R_2^\top=-L(F_{12}).
\]
These sign conditions force $L(F_{12})$ to have support only in the
$(1,2)$ and $(2,1)$ entries, hence
\[
L(F_{12})=a_{12}F_{12}
\]
for some scalar $a_{12}$. By permutation equivariance, the same scalar occurs
for every $F_{ij}$, so there is $\alpha\in\bbR$ with
\[
L(F_{ij})=\alpha F_{ij}
\qquad\text{for all }i<j.
\]

Now consider $D_{12}$. Every diagonal sign flip fixes $D_{12}$, so $L(D_{12})$
is fixed by all diagonal sign conjugations. Therefore $L(D_{12})$ is diagonal.
Permutations of coordinates $3,\dots,d$ show that all diagonal entries outside
the first two are equal, so
\[
L(D_{12})=\operatorname{diag}(a,b,c,\dots,c).
\]
Let $P_{12}$ be the permutation matrix exchanging coordinates $1$ and $2$.
Because
\[
P_{12}D_{12}P_{12}^\top=-D_{12},
\]
equivariance gives
\[
P_{12}L(D_{12})P_{12}^\top=-L(D_{12}).
\]
Thus $b=-a$ and $c=0$, hence
\[
L(D_{12})=aD_{12}.
\]
Let $Q_{\pi/4}$ be the rotation by $\pi/4$ in the $(1,2)$-plane. A direct
calculation gives
\[
Q_{\pi/4}D_{12}Q_{\pi/4}^\top=F_{12}.
\]
Equivariance then yields
\[
L(F_{12})
=
L(Q_{\pi/4}D_{12}Q_{\pi/4}^\top)
=
Q_{\pi/4}L(D_{12})Q_{\pi/4}^\top
=
aF_{12}.
\]
Since already $L(F_{12})=\alpha F_{12}$, one has $a=\alpha$. By permutation
equivariance,
\[
L(D_{ij})=\alpha D_{ij}
\qquad\text{for all }i<j.
\]

The matrices $\{F_{ij}\}_{i<j}$ together with the differences
$\{D_{ij}\}_{i<j}$ span the traceless symmetric carrier
$\mathrm{Sym}^2_0(\bbR^d)$. Hence
\[
L(S_0)=\alpha S_0
\qquad\text{for every }S_0\in \mathrm{Sym}^2_0(\bbR^d).
\]
Every symmetric tensor decomposes uniquely as
\[
S=\Bigl(S-\frac{\operatorname{tr}(S)}{d}I\Bigr)
+\frac{\operatorname{tr}(S)}{d}I
\]
with traceless part and trace part. Using the formulas above gives
\[
L(S)
=
\alpha\Bigl(S-\frac{\operatorname{tr}(S)}{d}I\Bigr)
+
\beta\,\frac{\operatorname{tr}(S)}{d}I.
\]
Uniqueness of $\alpha$ and $\beta$ follows from the direct-sum decomposition
\[
\mathrm{Sym}^2(\bbR^d)=\mathrm{Sym}^2_0(\bbR^d)\oplus \bbR I.
\qedhere
\]
\end{proof}

\begin{proof}[Proof of Theorem~\ref{thm:symmetry-reduction-admissibility}]
\label{proof:thm:symmetry-reduction-admissibility}
(i) Equivariant operators preserve isotypic components: $P_\rho A P_{\rho'} = 0$ for $\rho \ne \rho'$ when $A \in \mathrm{End}_G(\cH_{\mathrm{obs}})$.

(ii) For $x = \sum_\rho x_\rho$ with $x_\rho \in \cH_\rho$,
\[
\langle x, (D + Q) x \rangle = \sum_\rho \langle x_\rho, (D_\rho + Q_\rho) x_\rho \rangle.
\]
This is nonnegative for all $x$ iff each block summand is nonnegative.

(iii) Suppose $Q' \in \mathrm{End}_G(\cH_{\mathrm{obs}})$ satisfies $Q' \preceq Q$ with $D + Q' \succeq 0$.
Then $Q'_\rho \preceq Q_\rho$ and $D_\rho + Q'_\rho \succeq 0$ for each $\rho$.
Minimality of $Q_\rho$ forces $Q'_\rho = Q_\rho$, hence $Q' = Q$.
The converse assembles blockwise minima into a global minimum.

(iv) If $\langle x_*, (D + Q) x_* \rangle = 0$ with $x_* = \sum_\rho x_{*,\rho}$, then each summand $\langle x_{*,\rho}, (D_\rho + Q_\rho) x_{*,\rho} \rangle = 0$ (nonnegative terms summing to zero).
Any nonzero $x_{*,\rho}$ is a tight mode.
Minimal faces factor because the PSD cone in $\mathrm{End}_G(\cH_{\mathrm{obs}})$ is the product of PSD cones in each $\mathrm{End}_G(\cH_\rho)$.
\end{proof}

\begin{proof}[Proof of Theorem~\ref{thm:minimal-admissible-transport-order}]
\label{proof:thm:minimal-admissible-transport-order}
Because every admissible candidate has finite transport order, the set
\[
\{\operatorname{ord}_{\mathcal T}(A):A\in\mathcal F_{\mathrm{adm}}\}
\subseteq \bbN
\]
is nonempty. Every nonempty subset of $\bbN$ has a least element, which is
exactly $m_\ast$, proving (i). Statement (ii) is then just the definition of
the minimal-order admissible subclass. Statement (iii) follows because once the
minimal admissible order class has only one equivalence class, no second
inequivalent admissible candidate of the same minimal order remains.
\end{proof}

\begin{proof}[Proof of Theorem~\ref{thm:transport-generation-forces-scalar-recursion}]
\label{proof:thm:transport-generation-forces-scalar-recursion}
By item (i) of
Definition~\ref{def:transport-generated-admissible-family},
Theorem~\ref{thm:order-one-forcing-criterion} gives
\[
m_\ast=1.
\]
By item (iii) of
Definition~\ref{def:transport-generated-admissible-family},
Corollary~\ref{cor:no-preferred-direction-quadratic-transport} gives a unique
scalar $\kappa(B)\ge 0$ such that
\[
B^*B=\kappa(B)\,I_F.
\]
Finally, Theorem~\ref{thm:reversible-transport-recursion} gives the recursion
identity
\[
\mathcal R_Bx=0
\qquad\Longleftrightarrow\qquad
x_{n+1}=Bx_n-x_{n-1}
\]
for every bi-infinite sequence propagated by $B$. This proves all three
statements.
\end{proof}

% ============================================================
% Proofs from admissibility.tex
% ============================================================

\begin{proof}[Proof of Theorem~\ref{thm:scalar-completion}]
\label{proof:thm:scalar-completion}
(i) The condition $D + \nu L \succeq 0$ means $\langle x, (D + \nu L) x \rangle \ge 0$ for all $x$.
For $x \in \Ker(L)$, this reduces to $\langle x, D x \rangle \ge 0$, which holds by kernel compatibility.
For $x \notin \Ker(L)$, we have $\langle x, L x \rangle > 0$, so
\[
\langle x, (D + \nu L) x \rangle \ge 0
\quad\Leftrightarrow\quad
\nu \ge \frac{-\langle x, D x \rangle}{\langle x, L x \rangle}.
\]
Taking the supremum over all such $x$ gives the threshold $\nu_*$.

(ii) Since $\nu \mapsto \nu L$ is monotone in PSD order, $\nu_*$ is the unique minimal $\nu$ satisfying (i).

(iii) The maximizer $x_*$ achieves equality in the Rayleigh quotient, so $\langle x_*, (D + \nu_* L) x_* \rangle = 0$.
\end{proof}
